\chapter{The Torsiton Couplings from the Hehl--Datta Scale}
\label{app:couplings}

Chapter~\ref{chap:solitons} writes the torsiton as a bound state of one nonlinear Dirac field, but
the self-consistent solver needs numbers: the scalar coupling $g$, the repulsive vector coupling
$g_v$, the meson masses $m_\sigma,m_\omega$, and the chiral-potential quartic $\lambda$. None of
these is free. They all descend from the single Hehl--Datta coupling $G_\mathrm{HD}=\tfrac34\kappa$ by
Nambu--Jona-Lasinio bosonisation \citep{nambu1961}, leaving exactly \emph{one} dimensionful input,
the cutoff $\Lambda$ (the scale of Chapter~\ref{chap:weltformel} and lattice target~L5). Everything
else is a ratio. This appendix derives them; the algebra is reproduced symbolically in
\texttt{sims/.../coupling\_derivation.py} (SymPy) and supplied numerically by
\texttt{sims/.../njl\_couplings.py}.

\section{The Fierz channels}
The interaction is a single axial--axial term, $\mathcal{L}_\mathrm{int}=-G_\mathrm{HD}\,(\bar\psi
\gamma_5\gamma^\mu\psi)(\bar\psi\gamma_5\gamma_\mu\psi)$. A Fierz transformation re-expresses it in
the bilinear channels (Section~\ref{sec:fierz}); carried out with the Dirac matrices it gives
\[
  c_S=\tfrac14,\quad c_P=-\tfrac14,\quad c_V=\tfrac12,\quad c_A=\tfrac12,\quad c_T=0,
\]
so the two mean-field-relevant couplings are the scalar $G_S=\tfrac14 G_\mathrm{HD}$ (attractive, the
chiral condensate) and the vector $G_V=\tfrac12 G_\mathrm{HD}$ (repulsive), with the
\emph{walking-favourable} equality $G_A=G_V$. The ratio that matters is
\begin{equation}
  G_V/G_S = 2 .
  \label{eq:GVGS}
\end{equation}

\section{Bosonisation and the loop}
Introduce auxiliary fields $\sigma$ (scalar) and $\omega_\mu$ (vector) and integrate out the
fermions, whose dynamical mass is $M$. The one-loop integrals (4D Euclidean, sharp cutoff
$\Lambda$, $x\equiv\Lambda^2/M^2$) are
\begin{equation}
  I_0=\frac{1}{16\pi^2}\Big[\Lambda^2-M^2\ln(1+x)\Big],\qquad
  I_1=\frac{1}{16\pi^2}\Big[\ln(1+x)-\frac{x}{1+x}\Big].
  \label{eq:loops}
\end{equation}
$I_0$ carries the quadratic divergence, $I_1$ the logarithmic one. The chiral condensate is the
filled Dirac sea,
\begin{equation}
  \langle\bar\psi\psi\rangle = -4 N_c M\,I_0 ,
  \label{eq:condensate}
\end{equation}
and the gap equation $M=-2G_S\langle\bar\psi\psi\rangle$ fixes $M$ against $\Lambda$ and $G_S$---the
one place the absolute scale enters. The pion decay constant is
\begin{equation}
  f_\pi^2 = 4 N_c M^2 I_1 ,
  \label{eq:fpi}
\end{equation}
and we define the scalar Yukawa coupling by the Goldberger--Treiman relation $g\equiv M/f_\pi$, so
$v=f_\pi=M/g$.

\section{The scalar sector}
Chiral NJL gives the $\sigma$ at the $\bar q q$ threshold,
\begin{equation}
  m_\sigma = 2M .
  \label{eq:msigma}
\end{equation}
Matching the double-well $U(\sigma)=\tfrac{\lambda}{4}(\sigma^2-v^2)^2$, whose curvature is
$m_\sigma^2=2\lambda v^2$, fixes the quartic:
\begin{equation}
  \boxed{\ \lambda=\frac{m_\sigma^2}{2v^2}=\frac{4M^2}{2v^2}=2g^2\ }
  \label{eq:lambda}
\end{equation}
---derived, not the value-$8$ placeholder the solver had been using.

\section{The vector sector}
The NJL vector pole gives $m_\omega=\sqrt6\,M$ (with $M\simeq330$~MeV this is $\simeq810$~MeV, the
$\rho$), and the KSRF relation $m_\omega^2=2 g_v^2 f_\pi^2$ fixes the repulsive coupling:
\begin{equation}
  \boxed{\ m_\omega=\sqrt6\,M,\qquad g_v=\sqrt{\frac{m_\omega^2}{2v^2}}=\sqrt3\,g\ }
  \label{eq:vector}
\end{equation}

\section{Consistency, and the resolution of the unbinding}
The decisive check is that the bosonised masses and couplings reproduce
Eq.~\eqref{eq:GVGS}:
\begin{equation}
  \frac{G_V}{G_S}=\frac{g_v^2/m_\omega^2}{g^2/m_\sigma^2}
  =\frac{3g^2/6M^2}{g^2/4M^2}=2\ \checkmark
  \label{eq:consistency}
\end{equation}
This is why the soliton unbound when the repulsion was first switched on: the naive guess
$g_v=g\sqrt2$ with $m_\omega=gv$ violates Eq.~\eqref{eq:consistency}. The true coupling is
\emph{stronger} ($\sqrt3\,g$) but much \emph{shorter-ranged} ($\sqrt6\,M\gg m_\sigma$); the two
recombine to the Fierz ratio $2$ in integrated strength, leaving the spatial \emph{shape}---and the
relativistic $G^2\!\mp\!F^2$ balance of the first-order Dirac solver---to decide binding. At the
derived, QCD-like couplings ($g\simeq3$) the leading mean-field bag is \emph{marginal}: it binds
only when the colour levels are filled, and shallowly. The missing depth is the Dirac sea's own
polarisation, Eq.~\eqref{eq:condensate}, computed not in the uniform vacuum but in the soliton
background---the chiral-quark-soliton mode sum, whose vacuum-subtracted condensate reinforces the
bag. Carrying that sum out, de-doubled and regularised, is the calculation that decides whether the
torsiton stands self-bound at its own couplings or hands the spectrum to the lattice (L4).

\section{The hedgehog, and the torsiton mass}
\label{sec:hedgehog-mass}
The single-channel bag was marginal; the chiral soliton binds the way the QCD nucleon does---through
the \emph{hedgehog}. The chiral field is rotated radially in isospin, $M_\mathrm{vac}\exp(i\gamma_5\,
\boldsymbol\tau\!\cdot\!\hat{\mathbf r}\,\theta(r))$, locking spin and isospin into the conserved
grand spin $K=\mathbf J+\mathbf T$. In the $K{=}0$ sector a quark level dives out of the continuum
deep into the mass gap as the chiral angle winds from $\theta(0)=\pi$ to $\theta(\infty)=0$ (one unit
of winding, baryon number one): the valence quark, bound (\texttt{cartasis\_sims.hedgehog}).

Minimising the leading-order energy of the unit-winding soliton over its size $r_0$,
\begin{equation}
  E(r_0) = N_c\,E_\mathrm{val}(r_0) + \underbrace{\int 4\pi r^2\,dr\;\tfrac12 f_\pi^2
  \Big[\theta'^2 + \tfrac{2\sin^2\theta}{r^2}\Big]}_{E_2(r_0)\ \propto\ r_0},
  \label{eq:hhE}
\end{equation}
exposes the same Derrick balance that governs every three-dimensional chiral soliton. At the derived
coupling $f_\pi=M/g$ with $g\simeq3$ (so $f_\pi\simeq0.33\,M$) the two-derivative term $E_2$ grows
\emph{linearly} in $r_0$, while the valence level dives only to $E_\mathrm{val}\simeq+0.2\,M$ at
$r_0\simeq1.5\,M^{-1}$ and does not cross zero until $r_0\simeq2\,M^{-1}$, where $E_2\simeq8\,M$.
There is therefore \emph{no} self-bound minimum below the constituent sum: the infimum is the trivial
limit $r_0\!\to\!0$, where $E_2\!\to\!0$ and the $N_c$ valence quarks unbind to the gap edge, i.e.
$E\to N_cM$, approached from above. This is not a failure of the framework; it is the textbook
result. The two-derivative term alone never binds a nucleon below $N_cM$. The question is whether the
higher orders of the same quark loop---the four-derivative (Skyrme) term, equivalently the full
regularised Dirac-sea Casimir energy---bend the curve into a stable interior minimum. Computing that
sea energy in full (next section) settles it: at these couplings they do \emph{not}, and the soliton
sits at the constituent sum.

The torsiton mass is therefore fixed near the constituent sum,
\begin{equation}
  \boxed{\ M_\mathrm{torsiton}\ \simeq\ N_c\,M(\Lambda)\ }
  \label{eq:torsitonmass}
\end{equation}
with $M(\Lambda)$ set by the gap equation~\eqref{eq:condensate} from the single Hehl--Datta cutoff
$\Lambda$---one dimensionful number for the whole spectrum. With $N_c=3$ this is $M_\mathrm{torsiton}
\simeq3\,M$, the analogue of the proton sitting at three constituent quarks. The precise binding
fraction (whether the sea cancellation leaves the mass a few per cent above or below $3M$) is the
delicate, regularisation-sensitive piece that the non-perturbative lattice---target~L4---settles; the
mean-field calculation here fixes the \emph{scale} and hands off the last few per cent. The
minimisation and the Derrick collapse are reproduced in \texttt{cartasis\_sims.hedgehog}
(\texttt{minimize\_b1\_soliton}).

\paragraph{The full sea sum.} Replacing the two-derivative estimate by the full regularised Dirac-sea
energy means diagonalising the hedgehog Hamiltonian in every grand-spin sector $K=\mathbf J+\mathbf T$
and summing the vacuum-subtracted negative-energy levels with the proper-time cutoff. The apparatus is
built in two stages. \texttt{cartasis\_sims.chiral\_quark\_soliton} provides the grand-spin angular
couplings (the matrix elements of $\boldsymbol\tau\!\cdot\!\hat{\mathbf r}$, verified by $(\boldsymbol
\tau\!\cdot\!\hat{\mathbf r})^2=1$) and the Dirac Hamiltonian, whose free spectrum is parity-degenerate
with no spurious in-gap level and whose $K{=}0$ sector reproduces the valence dive of
\texttt{hedgehog}. On a finite-difference radial grid, though, the high-$K$ centrifugal barrier is
poorly resolved: the grand-spin sum crawls ($\sim K^{-1.2}$) and the subtracted energy drifts with the
box. The cure---and the standard way CQSM is solved---is the Kahana--Ripka basis of free
spherical-Bessel Dirac partial waves in a box (\texttt{cartasis\_sims.kahana\_ripka}), where the
barrier and the $r^l$ origin behaviour are exact, $\boldsymbol\sigma\!\cdot\!\mathbf p$ acts through
Bessel recurrences, and the vacuum subtraction is basis-consistent. There the sum \emph{converges}
($\Delta\!\sim\!10^{-3}$ by $K_{\max}\!\sim\!16$) and is box-stable, and the small-amplitude limit
returns a clean gradient term with $f_\pi\simeq0.39\,M$, consistent with Eq.~\eqref{eq:fpi}.

With the sea converged, the B=1 mass $M_\mathrm{sol}=N_c\,\varepsilon_\mathrm{val}+E_\mathrm{sea}$ can be
minimised. The outcome confirms Eq.~\eqref{eq:torsitonmass} from an actual mode sum: $E_\mathrm{sea}$ is
positive and grows with the soliton size while the valence level dives only to $\varepsilon_\mathrm{val}
\simeq+0.2\,M$, so at the derived (weak, $f_\pi/M\simeq0.4$) couplings the soliton has no deep interior
minimum---it relaxes toward the point limit where $E_\mathrm{sea}\to0$ and the $N_c$ valence quarks sit
at the gap edge, i.e.\ the constituent sum. Across cutoffs $\Lambda=2.5\text{--}4\,M$ the infimum is
robustly $N_cM$, approached from above (e.g.\ $M_\mathrm{sol}\simeq3.6\,M$ at the most compact size for
$\Lambda=2.5\,M$). So, computed and converged, $M_\mathrm{torsiton}\simeq N_cM\simeq3\,M$. Whether the
\emph{full} self-consistent profile (beyond the one-parameter family minimised here) dips a few per
cent below $3M$ is $\Lambda$-sensitive and is lattice target~L4; the excited in-gap orbitals
(\texttt{ingap\_levels} at higher $K$) are the candidate higher torsiton generations.

\bigskip
\noindent\emph{Reproducibility.} The Fierz coefficients are in \texttt{cartasis\_sims.fierz}; the
loop integrals, condensate, and the relations \eqref{eq:lambda}, \eqref{eq:vector},
\eqref{eq:consistency} are derived symbolically in \texttt{cartasis\_sims.coupling\_derivation}
and returned numerically by \texttt{cartasis\_sims.njl\_couplings}.
